\documentclass[12pt,a4paper]{article}
% Preambule commun aux comptes rendus CR_*.tex.

\usepackage{array}
\newcolumntype{C}[1]{>{\centering\arraybackslash}m{#1}}
\usepackage{multirow}
\usepackage{geometry}
\geometry{a4paper, margin=1in}
\usepackage{graphicx}
\usepackage{float}
\usepackage{tikz}
\usetikzlibrary{arrows.meta,positioning}
\usepackage{pgfgantt}
\usepackage{pdflscape}
\usepackage{enumitem}
\usepackage{booktabs}
\usepackage{longtable}
\usepackage{ragged2e}
\usepackage{tabularx}
\usepackage{xparse}
\usepackage{ltablex}
\usepackage{xltabular}
\usepackage{subcaption}

\newcolumntype{Y}{>{\RaggedRight\arraybackslash}X}

\newcommand{\StyledTableSetup}{%
  \footnotesize
  \renewcommand{\arraystretch}{1.2}%
  \setlength{\tabcolsep}{4pt}%
}


% ============================================================
% IHEDN COVER TEMPLATE (XeLaTeX)
% ============================================================
% Compilation (from project root):
%   latexmk -pdf main.tex
%   LATEXMK_SHELL_ESCAPE=0 latexmk -pdf main.tex
%
% Optional environment controls from latexmkrc:
%   LATEXMK_ENGINE=xelatex        (default/enforced)
%   LATEXMK_SHELL_ESCAPE=1|0      (default 1, needed for SVG conversion)
%
% Required package for this template:
%   \usepackage{ihedn-cover}
%
% Note: ihedn-cover already loads needed internal packages
% (fontspec, graphicx, tikz, ragged2e, optional svg).
% ============================================================

% ----------------------------------------------------------------
% Package options (documented)
% ----------------------------------------------------------------
% Geometry scope:
%   apply-geometry            -> force A4 + zero margins while cover is built
%   no-apply-geometry         -> do not alter host geometry (default)
%
% Typesetting freeze:
%   global-typesetting        -> apply freeze globally (intrusive)
%   no-global-typesetting     -> no global freeze (default)
%
% Small-caps handling:
%   local-smallcaps-fix       -> \textsc -> \MakeUppercase only inside cover (default)
%   no-local-smallcaps-fix    -> disable local fix
%   global-smallcaps-fix      -> global override (intrusive)
%   no-global-smallcaps-fix   -> disable global fix (default)
%
% Main font handling:
%   global-mainfont           -> set Times New Roman globally (default)
%   local-mainfont            -> set Times only during cover rendering
%   no-mainfont-override      -> keep host document main font
%
% SVG support:
%   svg-support               -> enable SVG logos (default, needs shell-escape)
%   no-svg-support            -> disable SVG path handling (pdf/png only)
%
% Debug overlay:
%   debug                     -> show mm grid + bounding boxes on cover
%   no-debug                  -> hide debug overlay (default)
% ----------------------------------------------------------------
\usepackage[
  no-apply-geometry,
  no-global-typesetting,
  global-smallcaps-fix,
  local-smallcaps-fix,
  global-mainfont,
  svg-support,
  no-debug
]{ihedn-cover}

% ----------------------------------------------------------------
% Bibliography stack (BibLaTeX/Biber, style from subtree vendor/)
% ----------------------------------------------------------------
\usepackage{polyglossia}
\setdefaultlanguage{french}
\makeatletter
\g@addto@macro\captionsfrench{%
  \renewcommand{\tablename}{Tableau}%
  \renewcommand{\figurename}{Figure}%
}
\makeatother
\usepackage{csquotes}
\usepackage{xurl}
\usepackage{hyperref}
\makeatletter
\let\ihedn@orig@input@path\input@path
\def\input@path{{vendor/biblatex-sciences-po-fr/style/}}
\makeatother
\usepackage[
  backend=biber,
  style=sciencespo,
  hyperref=true,
  autolang=other
]{biblatex}
\AtBeginBibliography{\sloppy\setlength{\emergencystretch}{3em}}
\makeatletter
\let\input@path\ihedn@orig@input@path
\makeatother
\addbibresource{bibliographies/bib/rapport.bib}

% ----------------------------------------------------------------
% tcolorbox
% ----------------------------------------------------------------
\usepackage[most]{tcolorbox}
\tcbuselibrary{breakable}

% Example: use Arial for the full report + cover.
% Uncomment and switch package option to no-mainfont-override.
% \setmainfont{Arial}[
%   UprightFont    = *,
%   BoldFont       = * Bold,
%   ItalicFont     = * Italic,
%   BoldItalicFont = * Bold Italic
% ]

%==================================================
% Liste des propositions
%==================================================

\makeatletter

\newcommand{\listpropositionname}{Liste des propositions}

\newcommand{\listofpropositions}{
  \section*{\listpropositionname}
  \addcontentsline{toc}{section}{\listpropositionname}
  \@starttoc{lop}
}

\newcommand*\l@proposition{\@dottedtocline{1}{1.5em}{2.3em}}

\makeatother

%==================================================
% Compteur des propositions
%==================================================

\newcounter{proposition}

%==================================================
% Style graphique des propositions
%==================================================

\newtcolorbox{propositionbox}[2][]{
  lower separated=false,
  colback=white!80!gray,
  colframe=white!20!black,
  fonttitle=\bfseries,
  colbacktitle=white!30!gray,
  coltitle=black,
  enhanced,
  sharp corners,
  attach boxed title to top left={
    xshift=0.5cm,
    yshift=-2mm
  },
  title=#2,
  #1
}

%==================================================
% Environnement proposition
%==================================================

\newenvironment{proposition}[1]
{
  \refstepcounter{proposition}

  \addcontentsline{lop}{proposition}{%
    \protect\numberline{\theproposition}#1%
  }

  \begin{propositionbox}
  {Proposition~\theproposition~: #1}
}
{
  \end{propositionbox}
}

% Renommage Liste des figures en Table des figures
\addto\captionsfrench{%
  \renewcommand{\listfigurename}{Table des figures}
}

\newcommand{\heure}[2]{{\NoAutoSpacing #1:#2}}

\DeclareFieldFormat[misc]{title}{#1}

% Activer ce package avec l'option none retire la césure.
%\usepackage[none]{hyphenat}

\usepackage{adjustbox}
\usepackage{pdfpages}

\makeatletter
\newcommand{\IhednSetPdfPageCount}[2]{%
  \begingroup
  \edef\ihedn@pdfpath{#2}%
  \ifXeTeX
    \global#1=\XeTeXpdfpagecount "\ihedn@pdfpath" %
  \else
    \ifLuaTeX
      \saveimageresource{\ihedn@pdfpath}%
      \global#1=\lastsavedimageresourcepages\relax
    \else
      \pdfximage{\ihedn@pdfpath}%
      \global#1=\pdflastximagepages\relax
    \fi
  \fi
  \endgroup
}
\makeatother

\newcount\IhednIncludedPdfPageCount
\newcommand{\IhednRenderPdfPagesInBoxes}[1]{%
  \IhednSetPdfPageCount{\IhednIncludedPdfPageCount}{#1}%
  \foreach \IhednIncludedPdfPage in {1,...,\the\IhednIncludedPdfPageCount}{%
    \begin{tcolorbox}[
      colback=gray!5,
      colframe=black,
      boxrule=0.5pt,
      arc=3mm
    ]
    \centering
    \includegraphics[page=\IhednIncludedPdfPage,width=0.95\textwidth]{#1}
    \end{tcolorbox}
    \ifnum\IhednIncludedPdfPage<\IhednIncludedPdfPageCount
      \clearpage
    \fi
  }%
}

\usepackage{tablefootnote}


\begin{document}

% ----------------------------------------------------------------
% Cover keys (all available keys are shown below)
% ----------------------------------------------------------------
% Header block (top-right), logo, title, report lines, subtitle, committee.
\PageDeGardeIHEDN{
    header_1={Union-IHEDN},
    header_2={Association des auditeurs IHEDN},
    header_3={région Paris / Île de France},
    date={le \today},
    logo_path={Figures/logos_IHEDN/Logo_AR_16_PARIS_IDF.jpg},
    title={Crises majeures : quel rôle pour les citoyens engagés et les associations IHEDN ?},
    title_max_lines=2,
    title_fontsize={26.000pt},
    title_baselineskip={28.667pt},
    reportline_1={Rapport du Comité d'étude de l'Association régionale des auditeurs},
    reportline_2={de l'IHEDN région Paris / Île de France (AR 16)},
    %subtitle={Compte rendu de la réunion du 23 avril 2026},
    subtitle={Proposition relative à la gestion et au pilotage du projet de mise en place du questionnaire},
    committee={Comité d'étude, d'octobre 2025 à juin 2026, composé de : Mme Valérie Arlé-Faivre, \textit{présidente} ; M. Aurélien Duvignac-Rosa, \textit{rapporteur} ; M. Pascal Bayot Duponchel ; M. Jean-Jacques Cleynen ; Mme Valérie Cowan ; M. Alain Fontan ; M. Pierre-Marie Giard ; Mme Muriel Joyeux ; M. Bruno Leray ; M. Yves-Henri Renhas ; M. Marc Roure ; Mme Isabelle de Segonzac ; M. Gérard Turck ; Mme Marie Danielle Vasquez-Duchêne.}
}

% \section{Éléments de rappel}

% La réunion du 23 Avril 2026 s'inscrit dans la continuité des réflexions engagées
% par le comité d'études autour du rôle potentiel des membres de l'AR~16 dans
% le cadre de crises majeures.

% La séance a porté sur l'organisation et la mise en place des jalons dans le cadre de la rédaction du questionnaire.

% --- Corps du document en chiffres arabes ---
\pagenumbering{Roman}

\tableofcontents

\clearpage

\addcontentsline{toc}{section}{Table des figures}
\listoffigures

\clearpage

\addcontentsline{toc}{section}{Liste des tableaux}
\listoftables

\clearpage

% --- Corps du document en chiffres arabes ---
\pagenumbering{arabic}

\section{Présentation}

Le présent document vise à exposer le calendrier prévisionnel et l'analyse des risques inhérents à la mise en œuvre d'un questionnaire à destination des membres de l'AR 16, susceptible d'être étendu à d'autres associations régionales.

Il s'appuie sur les réflexions menées dans le cadre de l'élaboration du rapport du comité relatif au rôle des citoyens engagés et des associations de l'IHEDN dans la gestion des crises majeures.

\section{Planification et structuration des tâches}

Dans la continuité des travaux du comité, qui ont mis en évidence la nécessité de disposer d'une connaissance objectivée des ressources humaines de l'AR 16, la mise en œuvre d'un questionnaire structurant constitue une étape clef. Cette démarche vise à dépasser les représentations partielles afin de construire une vision consolidée des compétences mobilisables en situation de crise.

Le présent plan d'action détaille les modalités opérationnelles de conception, de déploiement et d'exploitation de ce questionnaire, dans une logique de gestion de projet étendue jusqu'à septembre 2027.

\subsection{Planification et structuration des tâches}

La réalisation du questionnaire s'inscrit dans une démarche progressive, structurée autour de plusieurs phases complémentaires :

\begin{enumerate}
    \item \textbf{Phase de cadrage stratégique}
    \begin{itemize}
        \item Sélection de trois associations régionales pilotes afin de tester et valider le dispositif ;
        \item Réalisation d'une étude bibliographique approfondie (octobre 2025 - avril 2026) permettant de consolider le cadre conceptuel et méthodologique ;
        \item Obtention d'une validation institutionnelle formelle (Union IHEDN, Bureau, CODIR), condition préalable à tout déploiement.
    \end{itemize}

    \item \textbf{Phase de conception du questionnaire}
    \begin{itemize}
        \item Définition de l'objectif général, des thématiques et de la problématique centrale ;
        \item Élaboration d'une trame structurée et d'une grille d'exploitation des données ;
        \item Conception du contenu détaillé et d'une maquette du questionnaire ;
        \item Choix des outils techniques et du support de diffusion.
    \end{itemize}

    \item \textbf{Phase de mobilisation des ressources}
    \begin{itemize}
        \item Identification des compétences nécessaires (experts en sciences sociales, spécialistes RGPD, partenaires académiques) ;
        \item Structuration des partenariats (Union IHEDN, IRSEM, institutions académiques) ;
        \item Estimation des ressources et planification détaillée.
    \end{itemize}
\end{enumerate}

\newpage

\begin{enumerate}[resume]

    \item \textbf{Phase de lancement et de déploiement}
    \begin{itemize}
        \item Mise en place d'un comité de pilotage ;
        \item Communication interne et externe sur les objectifs du projet ;
        \item Préparation et diffusion de la campagne d'emailing ;
        \item Administration du questionnaire auprès des membres.
    \end{itemize}

    \item \textbf{Phase d'exploitation et de valorisation}
    \begin{itemize}
        \item Traitement et analyse des données collectées ;
        \item Production d'indicateurs et de cartographies de compétences ;
        \item Rédaction d'un rapport de travail et capitalisation des résultats ;
        \item Restitution finale prévue le 15 septembre 2027.
    \end{itemize}
\end{enumerate}

Cette structuration s'inscrit dans une logique de cycle complet de projet, allant de la conception à la valorisation, en cohérence avec les objectifs stratégiques du comité visant à renforcer la capacité d'analyse et de contribution de l'IHEDN dans la gestion des crises.

\clearpage

\subsection{Organisation des acteurs : une gouvernance en triptyque}

Un tel projet nécessaite la mise en place de processus et entités identifiées clairement de manière à garantir un suivi méthodologique optimal.

Pour cela, une gouvernance en tryptique est proposée. Celle-ci ce composerait d'abord d'un comité de pilotage afin d'assurer l'orientation stratégique, la validation des grandes étapes et le lien avec les instances décisionnelles. Ensuite un comité scientifique apporterait une validation scientifique concrète et une conformité juridique du dispositif du questionnaire. Enfin, la conception méthodologique et la mise en oeuvre du questionnaire serait réalisée par l'équipe projet composé de membres opérationnels du comité.

Une telle gouvernance, illustrée dans la figure \ref{fig:gouvernance-triptyque}, permettrait ainsi d'articuler expertise, légitimité institutionnelle et capacité opérationnelle, condition indispensable à la réussite d'une enquête de cette ampleur.

\begin{figure}[H]
\centering
\scalebox{1.5}{%
\begin{tikzpicture}[
    governance circle/.style={draw=black, line width=1.1pt, fill=gray!3},
    governance title/.style={font=\bfseries\footnotesize, align=center, text width=3.8cm},
    governance body/.style={font=\scriptsize, align=center, text width=3.6cm, anchor=north, inner sep=0pt},
    interaction arrow/.style={-{Stealth[length=2.4mm]}, line width=0.65pt},
    interaction label/.style={font=\tiny, align=center, fill=white, inner sep=1.2pt}
]
\def\governanceradius{2.05}
\coordinate (pilotage) at (0,2.15);
\coordinate (scientifique) at (-3.15,-1.25);
\coordinate (projet) at (3.15,-1.25);

\draw[governance circle] (pilotage) circle[radius=\governanceradius cm];
\draw[governance circle] (scientifique) circle[radius=\governanceradius cm];
\draw[governance circle] (projet) circle[radius=\governanceradius cm];

\draw[interaction arrow] (0.72,0.40) to[bend left=12]
    node[interaction label, pos=0.50, sloped, above] {Orientation\\Validation} (1.52,-0.05);
\draw[interaction arrow] (1.06,0.05) to[bend left=12]
    node[interaction label, pos=0.50, sloped, below] {Reporting} (0.26,0.50);

\draw[interaction arrow] (-0.72,0.40) to[bend right=12]
    node[interaction label, pos=0.50, sloped, above] {Cadrage\\Arbitrage} (-1.52,-0.05);
\draw[interaction arrow] (-1.06,0.05) to[bend right=12]
    node[interaction label, pos=0.50, sloped, below] {Avis\\Conformité} (-0.26,0.50);

\draw[interaction arrow] (-1.25,-1.05) to[bend left=10]
    node[interaction label, pos=0.50, above] {Appui méthodologique\\et RGPD} (1.25,-1.05);
\draw[interaction arrow] (1.25,-1.45) to[bend left=10]
    node[interaction label, pos=0.50, below] {Besoins terrain\\Retours opérationnels} (-1.25,-1.45);

\node[governance title] at (0,2.80) {Comité de pilotage};
\node[governance body] at (0,2.42) {%
    Présidents des trois\\régions pilotes\\[0.08cm]
    Présidence de l'Union IHEDN\\[0.08cm]
    Membres du comité d'étude
};

\node[governance title] at (-3.15,-0.60) {Comité scientifique};
\node[governance body] at (-3.15,-0.98) {%
    Expert en études\\et enquêtes sociologiques\\[0.08cm]
    Juriste RGPD
};

\node[governance title] at (3.15,-0.60) {Équipe projet};
\node[governance body] at (3.15,-0.98) {%
    Membres opérationnels\\du comité
};
\end{tikzpicture}
}
\caption{Organisation proposée de la gouvernance du projet de questionnaire.}
\label{fig:gouvernance-triptyque}
\end{figure}

Pour garantir l'aboutissement de ce projet, il est cependant nécessaire d'effectuer une analyse préalable des risques inhérents.

\clearpage

\subsection{Analyse des risques et facteurs de vigilance}

Le tableau \ref{tab:risques} présente les éventuels risques identifiés, en cohérence avec les retours d'expérience issus des travaux du comité.

\begin{table}[h]
\centering
\begin{tabular}{|p{3cm}|p{5cm}|p{2cm}|p{2cm}|}
\hline
\multicolumn{1}{|c|}{\textbf{Risque}} &
\multicolumn{1}{c|}{\textbf{Description}} &
\multicolumn{1}{c|}{\textbf{Probabilité}} &
\multicolumn{1}{c|}{\textbf{Impact}} \\
\hline
Coordination & Difficulté de coordination entre les trois associations régionales pilotes & Moyenne & Élevé \\
\hline
Interprétation & Mauvaise lecture ou instrumentalisation des résultats, pouvant générer inertie ou incompréhension & Faible & Élevé \\
\hline
Participation & Retards ou faibles taux de réponse, compromettant la représentativité des données & Moyenne & Élevé \\
\hline
Ressources & Inadéquation potentielle entre les compétences requises pour assurer la qualité de conception et l'appropriation du questionnaire, et les ressources effectivement mobilisées & Faible & Élevé \\
\hline
\end{tabular}
\caption{Cartographie des risques du projet.}
\label{tab:risques}
\end{table}

À ces risques s'ajoutent des points de vigilance déjà identifiés par le comité, notamment la disponibilité réelle des membres, la nécessité d'un socle commun de compétences et l'importance d'une communication claire et pédagogique.

\subsection{Facteurs clefs de succès}

La réussite du projet repose sur plusieurs conditions déterminantes :

\begin{itemize}
    \item Une validation institutionnelle claire et un portage politique fort ;
    \item Une conception méthodologique rigoureuse garantissant la crédibilité des résultats ;
    \item Une communication efficace favorisant l'adhésion des membres ;
    \item Une exploitation concrète des résultats permettant de renforcer la lisibilité et la légitimité de l'AR 16 auprès des autorités.
\end{itemize}

Au-delà de sa dimension technique, ce projet constitue un levier structurant pour l'association, en transformant un capital individuel de compétences en une ressource collective mobilisable au service de la résilience nationale.

\clearpage

\subsection{Calendrier prévisionnel de mise en œuvre du questionnaire}

La figure \ref{fig:gantt-questionnaire-ar16} présente l'orchestration des différentes étapes du projet, depuis leur planification jusqu'à l'obtention et l'exploitation des résultats du questionnaire en septembre 2027.

\begin{figure}[H]
\centering
\makeatletter
\providecommand{\ganttvrulelabel}[4][]{%
  \begingroup
  \ganttset{#1}%
  \gtt@tsstojulian{#4}{\gtt@vrule@slot}%
  \gtt@juliantotimeslot{\gtt@vrule@slot}{\gtt@vrule@slot}%
  \pgfmathsetmacro\y@upper{%
    \gtt@lasttitleline * \ganttvalueof{y unit title}%
  }%
  \pgfmathsetmacro\y@lower{%
    \gtt@lasttitleline * \ganttvalueof{y unit title}%
    + (\gtt@currentline - \gtt@lasttitleline - 1)%
    * \ganttvalueof{y unit chart}%
  }%
  \pgfmathsetmacro\x@mid{%
    (\gtt@vrule@slot - 1 + \ganttvalueof{vrule offset})%
    * \ganttvalueof{x unit}%
  }%
  \pgfmathsetmacro\y@label{\y@lower - #2}%
  \draw [/pgfgantt/vrule]
    (\x@mid pt, \y@upper pt) -- (\x@mid pt, \y@label pt)
    node [/pgfgantt/vrule label node] {#3};%
  \endgroup
}
\makeatother

\begin{ganttchart}[
    hgrid,
    x unit=0.0147cm,
    y unit chart=0.58cm,
    time slot format=isodate,
    title/.style={fill=gray!20, draw=gray!60},
    title label font=\bfseries\fontsize{6.4pt}{7pt}\selectfont,
    bar label font=\fontsize{6.8pt}{7.4pt}\selectfont,
    group label font=\bfseries\fontsize{6.8pt}{7.4pt}\selectfont,
    bar height=0.60,
    group height=0.38,
    group right shift=0,
    group left shift=0,
    group peaks tip position=0,
    group peaks height=0.18,
    group peaks width=0.2,
    bar/.style={fill=gray!35, draw=gray!70},
    group/.style={fill=gray!70, draw=gray!80},
    vrule/.style={dashed, line width=0.6pt, draw=black!70},
    vrule offset=.5,
    vrule label font=\fontsize{5pt}{5.6pt}\selectfont,
    vrule label node/.append style={anchor=north, rotate=45, align=left, inner sep=0.5pt, text=black!80},
    canvas/.style={draw=gray!50},
    y unit title=0.58cm,
    title height=1,
    title label anchor/.style={below=-1.6ex},
    bar label node/.append style={text width=4.2cm, align=right},
    group label node/.append style={text width=4.2cm, align=right}
]{2025-10-01}{2027-09-30}

\gantttitle{2025}{92}
\gantttitle{2026}{365}
\gantttitle{2027}{273} \\
\gantttitle{Oct}{31}
\gantttitle{Nov}{30}
\gantttitle{Déc}{31}
\gantttitle{Jan}{31}
\gantttitle{Fév}{28}
\gantttitle{Mar}{31}
\gantttitle{Avr}{30}
\gantttitle{Mai}{31}
\gantttitle{Juin}{30}
\gantttitle{Juil}{31}
\gantttitle{Août}{31}
\gantttitle{Sep}{30}
\gantttitle{Oct}{31}
\gantttitle{Nov}{30}
\gantttitle{Déc}{31}
\gantttitle{Jan}{31}
\gantttitle{Fév}{28}
\gantttitle{Mar}{31}
\gantttitle{Avr}{30}
\gantttitle{Mai}{31}
\gantttitle{Juin}{30}
\gantttitle{Juil}{31}
\gantttitle{Août}{31}
\gantttitle{Sep}{30} \\

\ganttgroup{Cadrage stratégique}{2025-10-01}{2026-06-25} \\
\ganttbar[bar/.append style={fill=blue!35}]{Étude bibliographique}{2025-10-01}{2026-04-30} \\
\ganttbar[bar/.append style={fill=blue!25}]{Sélection des trois AR pilotes}{2026-04-01}{2026-05-31} \\
\ganttbar[bar/.append style={fill=blue!45}]{Validation institutionnelle}{2026-03-30}{2026-05-25} \\

\ganttgroup{Conception du questionnaire}{2026-02-01}{2026-12-31} \\
\ganttbar[bar/.append style={fill=green!35}]{Objectifs, problématique et thèmes}{2026-02-01}{2026-12-31} \\
\ganttbar[bar/.append style={fill=green!25}]{Trame et grille d'exploitation}{2026-06-01}{2026-12-31} \\
\ganttbar[bar/.append style={fill=green!45}]{Maquette, support et outils}{2026-06-01}{2026-08-31} \\

\ganttgroup{Mobilisation et gouvernance}{2026-07-01}{2026-12-31} \\
\ganttbar[bar/.append style={fill=orange!35}]{Comité de pilotage}{2026-07-01}{2026-09-30} \\
\ganttbar[bar/.append style={fill=orange!25}]{Comité scientifique}{2026-07-01}{2026-09-30} \\
\ganttbar[bar/.append style={fill=orange!45}]{Ressources et partenariats}{2026-07-01}{2026-09-30} \\
\ganttbar[bar/.append style={fill=green!55}]{Validation RGPD et consentement}{2026-10-01}{2026-12-31} \\
\ganttbar[bar/.append style={fill=orange!55}]{Communication projet}{2026-10-01}{2026-12-31} \\
\ganttbar[bar/.append style={fill=orange!65}]{Préparation de l'emailing}{2026-10-01}{2026-12-31} \\

\ganttgroup{Lancement et déploiement}{2027-01-04}{2027-03-31} \\
\ganttbar[bar/.append style={fill=red!30}]{Lancement du questionnaire}{2027-01-04}{2027-02-28} \\
\ganttbar[bar/.append style={fill=red!40}]{Administration et relances}{2027-02-28}{2027-03-31} \\

\ganttgroup{Exploitation et valorisation}{2027-04-01}{2027-09-15} \\
\ganttbar[bar/.append style={fill=purple!25}]{Traitement des résultats}{2027-04-01}{2027-05-31} \\
\ganttbar[bar/.append style={fill=purple!35}]{Analyse et cartographie des compétences}{2027-04-01}{2027-06-30} \\
\ganttbar[bar/.append style={fill=purple!45}]{Rédaction du rapport de travail}{2027-06-01}{2027-08-31} \\
\ganttbar[bar/.append style={fill=purple!55}]{Capitalisation et restitution}{2027-08-01}{2027-09-15}

\ganttvrulelabel[vrule/.append style={draw=blue!70!black}, vrule label node/.append style={anchor=north east}]{10}{\shortstack[l]{15/05/2026\\Validation CODIR / Bureau}}{2026-05-15}
\ganttvrulelabel[vrule/.append style={draw=blue!70!black}, vrule label node/.append style={anchor=north west}]{34}{\shortstack[l]{25/06/2026\\Restitution du rapport de comité}}{2026-06-25}
\ganttvrulelabel[vrule/.append style={draw=red!70!black}, vrule label node/.append style={anchor=north west}]{34}{\shortstack[l]{31/03/2027\\Clôture de l'enquête}}{2027-03-31}
\ganttvrulelabel[vrule/.append style={draw=purple!70!black}, vrule label node/.append style={anchor=north east}]{10}{\shortstack[l]{15/09/2027\\Restitution finale}}{2027-09-15}

\end{ganttchart}

\caption{Calendrier prévisionnel de mise en œuvre du questionnaire de cartographie des compétences des membres de l'AR 16.}
\label{fig:gantt-questionnaire-ar16}
\end{figure}





\end{document}